%% GENERATED by `pipeline run results` - do not edit by hand.
\begin{table}[htbp]
\centering
\caption{Contribuição da própria partição, alfabeto de ação por unidade.}\label{tab:delta-star}
\label{tab:delta-star-unit}
\resizebox{\textwidth}{!}{%
\begin{tabular}{lrrrrr}
\toprule
família & top-1 real & top-1 embaralhado & $\Delta^{*}$ top-1 & $\Delta^{*}$ ajustado & $\Delta^{*}$ cobertura \\
\midrule
markov & 0,712 & 0,508 & \textbf{0,204} & 0,186 & 0,088 \\
gru & 0,715 & 0,496 & \textbf{0,219} & 0,202 & 0,118 \\
lstm & 0,705 & 0,471 & \textbf{0,234} & 0,216 & 0,121 \\
transformer & 0,713 & 0,488 & \textbf{0,225} & 0,207 & 0,154 \\
gnn & 0,625 & 0,428 & \textbf{0,197} & 0,179 & 0,073 \\
dfg & 0,589 & 0,296 & \textbf{0,293} & 0,275 & 0,095 \\
\bottomrule
\end{tabular}
}
\par\vspace{2pt}\footnotesize\raggedright $\Delta^{*}$ é o estimando: a partição real menos o seu controle embaralhado de forma equivalente, que tem o mesmo número de grupos, a mesma distribuição de tamanhos e a mesma população, com a pertença permutada. Ajustar um modelo por família muda o que esse modelo vê, de modo que uma partição desta forma desloca a acurácia por si só; apenas o excesso sobre o controle é atribuível ao agrupamento. $\Delta^{*}$ ajustado subtrai adicionalmente o modelo global agregado sob ambas as partições, cancelando um pequeno artefato de ponderação. $\Delta^{*}$ cobertura aplica o mesmo contraste a revocação de unidades $\times$ top-1 $/(1+$transições não vistas$)$. $\Delta^{*}$ é um critério de admissão, não uma ordenação: cresce com o número de famílias, porque especializar compensa mais quando as famílias são pequenas, de modo que partições de granularidade diferente não devem ser ordenadas por ele. Os alfabetos nunca são tabulados em conjunto.
\end{table}

\begin{table}[htbp]
\centering
\caption{Contribuição da própria partição, alfabeto de roteamento.}\label{tab:delta-star-itinerary}
\begin{tabular}{lrrrrr}
\toprule
família & top-1 real & top-1 embaralhado & $\Delta^{*}$ top-1 & $\Delta^{*}$ ajustado & $\Delta^{*}$ cobertura \\
\midrule
markov & 0,698 & 0,507 & \textbf{0,191} & 0,176 & 0,337 \\
gru & 0,705 & 0,485 & \textbf{0,220} & 0,205 & 0,339 \\
lstm & 0,690 & 0,448 & \textbf{0,242} & 0,227 & 0,290 \\
transformer & 0,711 & 0,499 & \textbf{0,212} & 0,198 & 0,306 \\
gnn & 0,477 & 0,304 & \textbf{0,173} & 0,158 & 0,174 \\
dfg & 0,460 & 0,261 & \textbf{0,200} & 0,185 & 0,191 \\
\bottomrule
\end{tabular}
\par\vspace{2pt}\footnotesize\raggedright $\Delta^{*}$ é o estimando: a partição real menos o seu controle embaralhado de forma equivalente, que tem o mesmo número de grupos, a mesma distribuição de tamanhos e a mesma população, com a pertença permutada. Ajustar um modelo por família muda o que esse modelo vê, de modo que uma partição desta forma desloca a acurácia por si só; apenas o excesso sobre o controle é atribuível ao agrupamento. $\Delta^{*}$ ajustado subtrai adicionalmente o modelo global agregado sob ambas as partições, cancelando um pequeno artefato de ponderação. $\Delta^{*}$ cobertura aplica o mesmo contraste a revocação de unidades $\times$ top-1 $/(1+$transições não vistas$)$. $\Delta^{*}$ é um critério de admissão, não uma ordenação: cresce com o número de famílias, porque especializar compensa mais quando as famílias são pequenas, de modo que partições de granularidade diferente não devem ser ordenadas por ele. Os alfabetos nunca são tabulados em conjunto.
\end{table}

\begin{table}[htbp]
\centering
\caption{Contribuição da própria partição, alfabeto de ação.}\label{tab:delta-star-procedure}
\begin{tabular}{lrrrr}
\toprule
família & top-1 real & top-1 embaralhado & $\Delta^{*}$ top-1 & $\Delta^{*}$ ajustado \\
\midrule
markov & 0,733 & 0,587 & \textbf{0,146} & 0,152 \\
gru & 0,755 & 0,601 & \textbf{0,153} & 0,160 \\
lstm & 0,754 & 0,599 & \textbf{0,155} & 0,161 \\
transformer & 0,751 & 0,593 & \textbf{0,158} & 0,165 \\
gnn & 0,644 & 0,515 & \textbf{0,129} & 0,136 \\
dfg & 0,644 & 0,517 & \textbf{0,126} & 0,133 \\
\bottomrule
\end{tabular}
\par\vspace{2pt}\footnotesize\raggedright $\Delta^{*}$ é o estimando: a partição real menos o seu controle embaralhado de forma equivalente, que tem o mesmo número de grupos, a mesma distribuição de tamanhos e a mesma população, com a pertença permutada. Ajustar um modelo por família muda o que esse modelo vê, de modo que uma partição desta forma desloca a acurácia por si só; apenas o excesso sobre o controle é atribuível ao agrupamento. $\Delta^{*}$ ajustado subtrai adicionalmente o modelo global agregado sob ambas as partições, cancelando um pequeno artefato de ponderação. $\Delta^{*}$ cobertura aplica o mesmo contraste a revocação de unidades $\times$ top-1 $/(1+$transições não vistas$)$. $\Delta^{*}$ é um critério de admissão, não uma ordenação: cresce com o número de famílias, porque especializar compensa mais quando as famílias são pequenas, de modo que partições de granularidade diferente não devem ser ordenadas por ele. Os alfabetos nunca são tabulados em conjunto.
\end{table}
